% GENERATED FILE. Edit canonical lessons, then run npm run generate:books.
% canonical-source-hash: fnv1a64:d58ef00c14cead10
% canonical-lessons: BN-C457-udar, BN-C457-kripon, BN-C457-antorik, BN-C457-poropokari, BN-C457-mohot

\chapter{Generous and Mean}
\label{ch:bn-generous-and-mean}
\hlchaptermodality{457}

\begin{chapteropening}
\textbf{By the end of this chapter:} \emph{I can say whether someone is generous or mean.}

Complete the last lesson of chapter 457: I can say whether someone is generous or mean.
\end{chapteropening}

\section[\texorpdfstring{\bn{উদার} (udār)}{উদার (udār)}]{\bn{উদার} (udār) — generous; broad-minded}

\label{lesson:BN-C457-udar}



\begin{quote}
Before the new one: say the Bengali for conduct, then the Bengali for to master.
\end{quote}

\subsection*{You'll want to know: \bn{উদার}}
\textbf{\bn{উদার}} — \emph{udār} — \textquotedblleft{}generous; broad-minded\textquotedblright{}.

\begin{quote}
\bn{আমার} \bn{কাকা} \bn{উদার}, \bn{সবসময়} \bn{গরিবদের} \bn{সাহায্য} \bn{করেন।}
\end{quote}

\begin{quote}
\emph{āmār kākā udār, sôbsômôy gôribder sāhājjô kôren.}
\end{quote}

\begin{quote}
My uncle is generous: he always helps the poor.
\end{quote}

\begin{cousinweb}
\textbf{\bn{উদার}} (\emph{udār}) is a tatsama, Sanskrit \emph{udāra}, \textquotedblleft{}high, lofty, noble\textquotedblright{} and \textquotedblleft{}generous\textquotedblright{} (Monier-Williams).
\end{cousinweb}

\begin{grammarlens}[title={What you've built}]
One more word for describing what you lived and felt.
\end{grammarlens}

\subsection*{Your turn}
\begin{itemize}
\raggedright
  \item \emph{Say it:} \emph{udār}
  \item \emph{Say it:} \emph{udār}, once more
  \item \emph{Say it:} \emph{āmār kākā udār, sôbsômôy gôribder sāhājjô kôren.} — My uncle is generous: he always helps the poor.
  \item \emph{Recall:} say the Bengali for conduct, then the Bengali for to master, then say \emph{udār} again
\end{itemize}

\begin{tcolorbox}[breakable,colback=teal!4,colframe=teal!35!black,title={Before you move on}]
What does \bn{উদার} mean? (\textquotedblleft{}Generous; broad-minded\textquotedblright{}.) Say it once more.
\end{tcolorbox}
\section[\texorpdfstring{\bn{কৃপণ} (kripôn)}{কৃপণ (kripôn)}]{\bn{কৃপণ} (kripôn) — miserly; a miser}

\label{lesson:BN-C457-kripon}



\begin{quote}
Before the new one: say the Bengali for to master, then the Bengali for generous.
\end{quote}

\subsection*{You'll want to know: \bn{কৃপণ}}
\textbf{\bn{কৃপণ}} — \emph{kripôn} — \textquotedblleft{}miserly; a miser\textquotedblright{}.

\begin{quote}
\bn{লোকটার} \bn{অনেক} \bn{টাকা} \bn{আছে}, \bn{কিন্তু} \bn{সে} \bn{খুব} \bn{কৃপণ।}
\end{quote}

\begin{quote}
\emph{lokṭār ônek ṭākā āchhe, kintu she khub kripôn.}
\end{quote}

\begin{quote}
The man has plenty of money, but he is a real miser.
\end{quote}

\begin{cousinweb}
\textbf{\bn{কৃপণ}} (\emph{kripôn}) is a tatsama, Sanskrit \emph{kṛpaṇa}, \textquotedblleft{}miserable, pitiable\textquotedblright{} (CDIAL 3436), which Monier-Williams also gives as \textquotedblleft{}niggardly\textquotedblright{}; that is the sense Bengali kept.
\end{cousinweb}

\begin{grammarlens}[title={What you've built}]
One more word for describing what you lived and felt.
\end{grammarlens}

\subsection*{Your turn}
\begin{itemize}
\raggedright
  \item \emph{Say it:} \emph{kripôn}
  \item \emph{Say it:} \emph{kripôn}, once more
  \item \emph{Say it:} \emph{lokṭār ônek ṭākā āchhe, kintu she khub kripôn.} — The man has plenty of money, but he is a real miser.
  \item \emph{Recall:} say the Bengali for to master, then the Bengali for generous, then say \emph{kripôn} again
\end{itemize}

\begin{tcolorbox}[breakable,colback=teal!4,colframe=teal!35!black,title={Before you move on}]
What does \bn{কৃপণ} mean? (\textquotedblleft{}Miserly; a miser\textquotedblright{}.) Say it once more.
\end{tcolorbox}
\section[\texorpdfstring{\bn{আন্তরিক} (āntôrik)}{আন্তরিক (āntôrik)}]{\bn{আন্তরিক} (āntôrik) — warm-hearted; sincere}

\label{lesson:BN-C457-antorik}



\begin{quote}
Before the new one: say the Bengali for generous, then the Bengali for miserly.
\end{quote}

\subsection*{You'll want to know: \bn{আন্তরিক}}
\textbf{\bn{আন্তরিক}} — \emph{āntôrik} — \textquotedblleft{}warm-hearted; sincere\textquotedblright{}.

\begin{quote}
\bn{আমাদের} \bn{নতুন} \bn{প্রতিবেশীরা} \bn{আন্তরিক।}
\end{quote}

\begin{quote}
\emph{āmāder nôtun prôtibeshīrā āntôrik.}
\end{quote}

\begin{quote}
Our new neighbours are warm-hearted.
\end{quote}

\begin{cousinweb}
\textbf{\bn{আন্তরিক}} (\emph{āntôrik}) is Sanskrit \emph{antara}, \textquotedblleft{}interior, intimate\textquotedblright{} (Monier-Williams), with the adjective ending \emph{-ika}: from the inside, from the heart.
\end{cousinweb}

\begin{grammarlens}[title={What you've built}]
One more word for describing what you lived and felt.
\end{grammarlens}

\subsection*{Your turn}
\begin{itemize}
\raggedright
  \item \emph{Say it:} \emph{āntôrik}
  \item \emph{Say it:} \emph{āntôrik}, once more
  \item \emph{Say it:} \emph{āmāder nôtun prôtibeshīrā āntôrik.} — Our new neighbours are warm-hearted.
  \item \emph{Recall:} say the Bengali for generous, then the Bengali for miserly, then say \emph{āntôrik} again
\end{itemize}

\begin{tcolorbox}[breakable,colback=teal!4,colframe=teal!35!black,title={Before you move on}]
What does \bn{আন্তরিক} mean? (\textquotedblleft{}Warm-hearted; sincere\textquotedblright{}.) Say it once more.
\end{tcolorbox}
\section[\texorpdfstring{\bn{পরোপকারী} (pôropôkārī)}{পরোপকারী (pôropôkārī)}]{\bn{পরোপকারী} (pôropôkārī) — helpful to others; kind-hearted}

\label{lesson:BN-C457-poropokari}



\begin{quote}
Before the new one: say the Bengali for miserly, then the Bengali for sincere.
\end{quote}

\subsection*{You'll want to know: \bn{পরোপকারী}}
\textbf{\bn{পরোপকারী}} — \emph{pôropôkārī} — \textquotedblleft{}helpful to others; kind-hearted\textquotedblright{}.

\begin{quote}
\bn{পরোপকারী} \bn{ছেলেটা} \bn{সবসময়} \bn{বুড়োদের} \bn{সাহায্য} \bn{করে।}
\end{quote}

\begin{quote}
\emph{pôropôkārī chheleṭā sôbsômôy buṛoder sāhājjô kôre.}
\end{quote}

\begin{quote}
The kind-hearted boy always helps old people.
\end{quote}

\begin{cousinweb}
\textbf{\bn{পরোপকারী}} (\emph{pôropôkārī}) is a tatsama, Sanskrit \emph{paropakārin}, \textquotedblleft{}assisting others, beneficent\textquotedblright{} (Monier-Williams): \emph{para} \textquotedblleft{}another\textquotedblright{} with \emph{upakāra} \textquotedblleft{}help\textquotedblright{}.
\end{cousinweb}

\begin{grammarlens}[title={What you've built}]
One more word for describing what you lived and felt.
\end{grammarlens}

\subsection*{Your turn}
\begin{itemize}
\raggedright
  \item \emph{Say it:} \emph{pôropôkārī}
  \item \emph{Say it:} \emph{pôropôkārī}, once more
  \item \emph{Say it:} \emph{pôropôkārī chheleṭā sôbsômôy buṛoder sāhājjô kôre.} — The kind-hearted boy always helps old people.
  \item \emph{Recall:} say the Bengali for miserly, then the Bengali for sincere, then say \emph{pôropôkārī} again
\end{itemize}

\begin{tcolorbox}[breakable,colback=teal!4,colframe=teal!35!black,title={Before you move on}]
What does \bn{পরোপকারী} mean? (\textquotedblleft{}Helpful to others; kind-hearted\textquotedblright{}.) Say it once more.
\end{tcolorbox}
\section[\texorpdfstring{\bn{মহৎ} (môhôt)}{মহৎ (môhôt)}]{\bn{মহৎ} (môhôt) — noble; great-hearted}

\label{lesson:BN-C457-mohot}



\begin{quote}
Before the new one: say the Bengali for sincere, then the Bengali for helpful to others.
\end{quote}

\subsection*{You'll want to know: \bn{মহৎ}}
\textbf{\bn{মহৎ}} — \emph{môhôt} — \textquotedblleft{}noble; great-hearted\textquotedblright{}.

\begin{quote}
\bn{গরিব} \bn{ছাত্রটির} \bn{পড়ার} \bn{খরচ} \bn{দিয়ে} \bn{তিনি} \bn{একটা} \bn{মহৎ} \bn{কাজ} \bn{করলেন।}
\end{quote}

\begin{quote}
\emph{gôrib chhātrôṭir pôṛār khôrôch diye tini ekṭā môhôt kāj kôrlen.}
\end{quote}

\begin{quote}
By paying for the poor student's studies, he did a noble thing.
\end{quote}

\begin{cousinweb}
\textbf{\bn{মহৎ}} (\emph{môhôt}) is a tatsama, Sanskrit \emph{mahat}, \textquotedblleft{}great, high, eminent\textquotedblright{} (Monier-Williams); Bengali uses it for greatness of heart. It ends in the new letter \bn{ৎ}.
\end{cousinweb}

\begin{grammarlens}[title={What you've built}]
One more word for describing what you lived and felt.
\end{grammarlens}

\subsection*{Your turn}
\begin{itemize}
\raggedright
  \item \emph{Say it:} \emph{môhôt}
  \item \emph{Say it:} \emph{môhôt}, once more
  \item \emph{Say it:} \emph{gôrib chhātrôṭir pôṛār khôrôch diye tini ekṭā môhôt kāj kôrlen.} — By paying for the poor student's studies, he did a noble thing.
  \item \emph{Recall:} say the Bengali for sincere, then the Bengali for helpful to others, then say \emph{môhôt} again
\end{itemize}

\begin{tcolorbox}[breakable,colback=teal!4,colframe=teal!35!black,title={Before you move on}]
What does \bn{মহৎ} mean? (\textquotedblleft{}Noble; great-hearted\textquotedblright{}.) Say it once more.
\end{tcolorbox}
